\documentclass[11pt,a4paper]{article}

% --- Essential Packages ---
\usepackage[utf8]{inputenc}
\usepackage{amsmath,amssymb,amsfonts,amsthm}
\newtheorem{prop}{Proposition}
\newtheorem{thm}{Theorem}
\usepackage{graphicx}
\usepackage{hyperref}
\usepackage{geometry}
\usepackage{cite}
\usepackage{microtype}
\usepackage{booktabs}
\usepackage{placeins}

% --- Page Geometry & Setup ---
\geometry{margin=1.0in}
\hypersetup{colorlinks=true, linkcolor=blue, citecolor=blue, urlcolor=blue}

\title{\Large\bfseries Substrate Cognitive Dynamics: Hierarchical Neural Networks, Predictive Processing, Memory Integration, and Emergent Cognition}

\author{
	\textbf{Ryan Beem}\\[0.5em]
	\small Independent Researcher \\
	\small \href{mailto:ryan@formoreoptions.com}{ryan@formoreoptions.com}
}

\date{\today}

\begin{document}
	
	\maketitle
	
	\begin{abstract}
		Building upon the biological process mechanics established in Paper X, we formulate a deterministic continuum framework for neural information processing, predictive processing, working memory, attention, decision making, and emergent cognition. In standard neuroscience, cognition is modeled as statistical rate-coding across discrete connectionist neurons. In the Unified Substrate Theory, we prove that the central nervous system is a multi-scale, coupled order-parameter phase network $\mathbf{N}_{\text{brain}}(\mathbf{x}, t) = \sum \mathbf{n}_i(\mathbf{x}, t) \in S^2$ operating as a global gradient optimizer. We derive the Predictive Error Minimization Theorem, demonstrating that internal generative models reduce sensory phase mismatch ($\Delta \Phi \neq 0$) by adjusting network streamtube topologies. We model attention as selective phase amplification ($A_0 + \Delta A$), formulate decision making as competitive energy relaxation across nested attractor basins, and re-frame collective neural oscillations (alpha, beta, gamma, theta) as global acoustic standing-wave modes.
	\end{abstract}
	
	\vspace{1em}
	\hrule
	\vspace{1.5em}
	
	% =========================================================
	% SECTION 1: FROM SINGLE NEURONS TO COUPLED PHASE NETWORKS
	% =========================================================
	\section{From Single Neurons to Massively Coupled Phase Networks}
	\label{sec:intro_cognitive_networks}
	
	In Paper X \cite{PaperX}, individual axonal action potentials were derived as solitary substrate phase pulses $\mathbf{n}(z - v_{\text{nerve}} t)$ propagating along lipid-sheathed channels. In cognitive neuroscience, intelligence emerges from the collective interaction of over $10^{11}$ interconnected neural conduits.
	
	In the Unified Substrate Theory (UST), the central nervous system is modeled as a continuous, multi-scale order-parameter phase field $\mathbf{N}_{\text{brain}}(\mathbf{x}, t) \in S^2$:
	\begin{equation}
		\mathbf{N}_{\text{brain}}(\mathbf{x}, t) = \sum_{i=1}^{N_{\text{neurons}}} \mathbf{n}_i(\mathbf{x}, t) e^{i \Phi_i(\mathbf{x}, t)}.
		\label{eq:brain_phase_field}
	\end{equation}
	
	The network evolves via coupled order-parameter strain relaxation:
	\begin{equation}
		\frac{\partial \mathbf{N}_{\text{brain}}}{\partial t} = -D_{\text{cortex}} \nabla^2 \mathbf{N}_{\text{brain}} - \gamma \frac{\delta E_{\text{network}}}{\delta \mathbf{N}_{\text{brain}}} + \mathbf{S}_{\text{sensory}}(\mathbf{x}, t),
		\label{eq:cortical_field_evolution}
	\end{equation}
	where $\mathbf{S}_{\text{sensory}}$ represents incoming order-parameter perturbations driven by peripheral sensory organs.
	
	% =========================================================
	% SECTION 2: PREDICTIVE ERROR MINIMIZATION THEOREM
	% =========================================================
	\section{Predictive Error Minimization and Internal Generative Models}
	\label{sec:predictive_processing}
	
	Rather than passively processing sensory inputs, the cortex actively generates top-down predictive phase patterns to anticipate environmental perturbations.
	
	\begin{thm}[Predictive Error Minimization Theorem]
		\label{thm:predictive_error}
		The total information-processing cost $E_{\text{prediction}}$ incurred by a cortical network is the spatial order-parameter strain differential between bottom-up sensory phase inputs $\nabla \mathbf{N}_{\text{sensory}}$ and top-down predictive model waves $\nabla \mathbf{N}_{\text{model}}$:
		\begin{equation}
			E_{\text{prediction}} = \frac{1}{2}\mu_0 \int_{\Omega_{\text{cortex}}} \left| \nabla \mathbf{N}_{\text{sensory}}(\mathbf{x}, t) - \nabla \mathbf{N}_{\text{model}}(\mathbf{x}, t) \right|^2 d^3x.
			\label{eq:predictive_error_functional}
		\end{equation}
		Learning adjusts synaptic streamtube topology (Paper X) along steepest-descent trajectories to minimize $E_{\text{prediction}}$ ($\Delta E_{\text{prediction}} \to 0$).
	\end{thm}
	
	\begin{proof}
		When top-down predictive waves match incoming sensory phase fronts ($\mathbf{N}_{\text{model}} \approx \mathbf{N}_{\text{sensory}}$), destructive strain interference eliminates localized phase gradients ($\nabla \mathbf{N}_{\text{sensory}} - \nabla \mathbf{N}_{\text{model}} = \mathbf{0}$). The net strain energy functional $E_{\text{prediction}}$ drops to a local minimum, suppressing prediction error signals and conserving metabolic ATP energy (Paper X). An unexpected sensory input produces a phase mismatch ($\Delta \Phi \neq 0 \implies \nabla \delta \mathbf{N} \gg 0$), propagating error pulses upward through the cortical hierarchy to update internal phase topology.
	\end{proof}
	
	% =========================================================
	% SECTION 3: WORKING MEMORY AS OSCILLATORY PHASE LOOPS
	% =========================================================
	\section{Working Memory as Persistent Oscillatory Phase Loops}
	\label{sec:working_memory}
	
	In Paper X, long-term memory was established as persistent structural phase-locking ($\Delta E_{\text{synapse}} < 0$). Working memory represents the short-term, active maintenance of information without permanent structural remodeling.
	
	\begin{prop}[Working Memory Phase Loop Invariant]
		\label{prop:working_memory_loop}
		Active working memory corresponds to a self-reinforcing, closed standing-wave phase loop traversing localized cortico-thalamic streamtube circuits. The active representation $M_{\text{active}}(t)$ is sustained by periodic phase recirculation:
		\begin{equation}
			M_{\text{active}}(t) = \oint_{C_{\text{loop}}} \mathbf{N}_{\text{brain}}(\mathbf{x}, t) \cdot d\mathbf{l} \equiv \text{const},
			\label{eq:working_memory_loop_integral}
		\end{equation}
		persisting until metabolic energy decay or external phase interference disrupts loop recirculation.
	\end{prop}
	
	% =========================================================
	% SECTION 4: ATTENTION AS SELECTIVE PHASE AMPLIFICATION
	% =========================================================
	\section{Attention as Selective Phase Amplification}
	\label{sec:attention_mechanics}
	
	Cognitive attention is the selective allocation of computational resources to specific sensory or internal channels.
	
	\begin{thm}[Selective Phase Amplification]
		\label{thm:attention_amplification}
		Top-down attentional modulation enhances the local signal-to-noise ratio of a target neural channel $\Omega_{\text{target}}$ by injectively scaling its phase amplitude vector from baseline $A_0$ to $A_0 + \Delta A$:
		\begin{equation}
			\mathbf{N}_{\text{attended}}(\mathbf{x}, t) = (A_0 + \Delta A) \, \mathbf{n}_{\text{base}}(\mathbf{x}, t) e^{i \Phi(\mathbf{x}, t)}.
			\label{eq:attention_scaling_formula}
		\end{equation}
	\end{thm}
	
	\begin{proof}
		Increasing local phase amplitude $\Delta A$ elevates the strain energy density $\mathcal{E}_{\text{target}} = \frac{1}{2}\mu_0 |\nabla \mathbf{N}_{\text{attended}}|^2$ relative to surrounding un-attended background channels ($\mathcal{E}_{\text{background}} \ll \mathcal{E}_{\text{target}}$). Attended phase fronts dominate global coupled network relaxation (Eq.\ \ref{eq:cortical_field_evolution}), forcing downstream decision attractors to align with the amplified target signal.
	\end{proof}
	
	% =========================================================
	% SECTION 5: DECISION FORMATION AS COMPETING ATTRACTOR BASINS
	% =========================================================
	\section{Decision Formation as Competitive Energy Basin Relaxation}
	\label{sec:decision_formation}
	
	In classical decision theory, choice selection is modeled via stochastic drift-diffusion processes. In UST, decision making is deterministic relaxation across competing topological energy basins.
	
	\begin{thm}[Competitive Attractor Decision Theorem]
		\label{thm:decision_attractor}
		Let $K$ mutually exclusive behavioral choices correspond to $K$ distinct local strain energy minima $\{E_1, E_2, \dots, E_K\}$ on the network state manifold $\mathcal{M}_{\text{cortex}}$. The network trajectory $\boldsymbol{\xi}(t)$ evolves deterministically toward the deepest accessible attractor basin:
		\begin{equation}
			\frac{d \boldsymbol{\xi}(t)}{dt} = -\nabla_{\boldsymbol{\xi}} E_{\text{network}}(\boldsymbol{\xi}) + \mathbf{S}_{\text{sensory}}(t).
			\label{eq:attractor_basin_flow}
		\end{equation}
		A commitment decision occurs when $\boldsymbol{\xi}(t)$ crosses a unstable phase threshold $\boldsymbol{\xi}^\ddagger$, collapsing the network state into a single stable motor action basin.
	\end{thm}
	
	% =========================================================
	% SECTION 6: COLLECTIVE NEURAL OSCILLATIONS (EEG BANDS)
	% =========================================================
	\section{Collective Neural Oscillations as Substrate Standing Waves}
	\label{sec:neural_oscillations}
	
	Electroencephalographic (EEG) frequency bands—delta ($1$--$4 \text{ Hz}$), theta ($4$--$8 \text{ Hz}$), alpha ($8$--$12 \text{ Hz}$), beta ($13$--$30 \text{ Hz}$), and gamma ($30$--$80 \text{ Hz}$)—are standard signatures of cognitive state.
	
	\begin{thm}[Substrate Neural Resonance Invariant]
		\label{thm:eeg_resonances}
		Macroscopic brain rhythms are the global acoustic standing-wave eigenmodes of the coupled order-parameter field $\mathbf{N}_{\text{brain}}$ bounded by cortical cavity geometry $\Omega_{\text{cortex}}$:
		\begin{equation}
			\nabla^2 (\delta \mathbf{N}) - \frac{1}{v_{\text{phase}}^2} \frac{\partial^2 (\delta \mathbf{N})}{\partial t^2} = \mathbf{0} \implies \omega_k = v_{\text{phase}} \cdot k,
			\label{eq:eeg_standing_wave_pde}
		\end{equation}
		where $v_{\text{phase}}$ is the effective intra-cortical shear conduction velocity (Paper VIII).
	\end{thm}
	
	\begin{proof}
		Gamma oscillations ($30$--$80 \text{ Hz}$) represent localized, short-range standing waves mediating fine-grained feature integration ($\nabla_{\text{local}} \mathbf{N}$). Theta ($4$--$8 \text{ Hz}$) and Alpha ($8$--$12 \text{ Hz}$) rhythms correspond to long-range, cross-cortical standing waves coordinating large-scale phase synchronization between sensory areas and sub-cortical memory structures (Paper X).
	\end{proof}
	
	% =========================================================
	% SECTION 7: THE COMPLETE ELEVEN-PAPER UST HIERARCHY
	% =========================================================
	\section{The Complete Eleven-Paper UST Scale Progression}
	\label{sec:eleven_paper_synthesis}
	
	With Paper XI, the Unified Substrate Theory constructs a continuous, 11-paper deterministic continuum spanning sub-atomic physics to emergent cognition:
	
	\begin{table}[htbp]
		\centering
		\caption{\textbf{The Unified Substrate Theory (UST) Complete Scale Ladder}}
		\begin{tabular}{lll}
			\toprule
			\textbf{Volume} & \textbf{Core Field Domain} & \textbf{Primary Governing Invariant} \\
			\midrule
			\textbf{Paper I} & Soliton Physics & Skyrme topological energy minimum \\
			\textbf{Paper II} & Hadronic Physics & Topological proton rest mass ($m_p$) \\
			\textbf{Paper III} & Electrodynamics & Fine-structure constant ($\alpha^{-1}$) \\
			\textbf{Paper IV} & Gravitation & Pressure-shadow deficit ($G$) \\
			\textbf{Paper V} & Atomic Chemistry & Orbital cavity eigenmodes \& $2n^2$ shells \\
			\textbf{Paper VI} & Organic Chemistry & Carbon $\sigma/\pi$ bonding \& Hückel $4n+2$ torus \\
			\textbf{Paper VII} & Materials Science & Streamtube network dimensionality ($D$) \\
			\textbf{Paper VIII} & Quantum Materials & Global macroscopic phase locking ($\nabla \Phi = \text{const}$) \\
			\textbf{Paper IX} & Self-Organization & Protein funnels \& DNA double-helix topology \\
			\textbf{Paper X} & Systems Biology & Metabolic strain reservoirs \& neural phase pulses \\
			\textbf{Paper XI} & Cognitive Dynamics & Predictive error minimization ($\min E_{\text{pred}}$) \\
			\bottomrule
		\end{tabular}
		\label{tab:ust_eleven_paper_summary}
	\end{table}
	
	% =========================================================
	% SECTION 8: CONCLUSION
	% =========================================================
	\section{Conclusion}
	\label{sec:conclusion}
	
	We have demonstrated that emergent cognition, predictive processing, working memory, attention, decision formation, and macro-scale EEG rhythms are direct consequences of coupled order-parameter phase-network optimization ($\min E_{\text{prediction}}$). By unifying physics, chemistry, biology, and cognitive science under a single continuum field framework, UST establishes that intelligence is a multi-scale, adaptive manifestation of substrate gradient-energy minimization.
	
	% =========================================================
	% REFERENCES
	% =========================================================
	\begin{thebibliography}{99}
		\bibitem{PaperX} R. Beem, \textit{Substrate Systems Biology: Metabolic Strain Reservoirs, Molecular Machinery, Neural Phase Pulses, and Emergent Adaptive Dynamics}, UST Paper X (2026).
		\bibitem{PaperIX} R. Beem, \textit{Substrate Molecular Self-Organization: Hydrogen Bonding, Protein Folding Topologies, DNA Helical Pairing, and Emergent Biological Structure}, UST Paper IX (2026).
		\bibitem{PaperVIII} R. Beem, \textit{Substrate Quantum Materials: Global Phase Coherence, Cooper-Pair Topological Analogs, and Ab-Initio Superconductivity}, UST Paper VIII (2026).
		\bibitem{PaperVII} R. Beem, \textit{Substrate Materials: Carbon Allotropes, Graphene Transport, Fracture Mechanics, and Condensed Matter Structure}, UST Paper VII (2026).
		\bibitem{PaperVI} R. Beem, \textit{Substrate Chemistry II: Carbon Bonding Topologies, Topological Aromaticity, and Deterministic Reaction Dynamics}, UST Paper VI (2026).
		\bibitem{PaperV} R. Beem, \textit{Substrate Chemistry I: Atomic Orbital Resonances, Topological Pauli Exclusion, Covalent Gradient Minimization, and Molecular Geometry}, UST Paper V (2026).
		\bibitem{PaperIV} R. Beem, \textit{Hydrostatic Pressure-Shadow Gravitation and $G$ Derivation}, UST Paper IV (2026).
		\bibitem{PaperIII} R. Beem, \textit{Substrate Electrodynamics: Fine-Structure Constant $\alpha$ and Precision Lepton Mass Hierarchy}, UST Paper III (2026).
		\bibitem{PaperII} R. Beem, \textit{Ab-Initio Derivation of Proton Mass, Structure, and Charge Radius}, UST Paper II (2026).
		\bibitem{PaperI} R. Beem, \textit{Topological Stabilization of 3D Solitons in Non-Linear Field Media}, UST Paper I (2026).
	\end{thebibliography}
	
\end{document}